\documentclass[english, parskip=full]{scrartcl}
\usepackage[utf8]{inputenc}  % Kodierung der Datei
\usepackage[english]{babel}  % Sprache auf Deutsch 
\usepackage{color}
\usepackage{graphicx, gensymb}
\usepackage{amsfonts, cancel, amssymb}
\usepackage{amsmath, amsthm, array, esint}
\usepackage{booktabs, multirow}  % schönere Tabellen
\usepackage{caption}  % Anpassung der Captions
\usepackage{scrlayer-scrpage}    % Manipulation des Seitenstils
\pagestyle{scrheadings}  % Stil für die Seite setzen
\automark{section}  % Abschnittsnamen als Seitenbeschriftung 
\setheadsepline{.5pt}    % Kopzeile durch Linie abtrennen
\usepackage[hidelinks]{hyperref}  % Links und weitere PDF-Features
\usepackage{typearea, tikz, pgffor, verbatim}
\usetikzlibrary{calc, arrows.meta,arrows, graphs, quotes, positioning}
\usepackage[cache=false, outputdir=build]{minted}
\usemintedstyle{vs}
\usepackage{placeins} % provides \FloatBarrier

\definecolor{bg}{rgb}{0.9,0.9,0.9}
\newcommand\numberthis{\addtocounter{equation}{1}\tag{\theequation}}
\usepackage{svg}
\usepackage{siunitx}
\usepackage{physics}
\usepackage{tensor}
\renewcommand{\bra}{}
\newcommand{\kom}{}
\newcommand{\brak}{}
\renewcommand{\braket}{}
\newcommand{\intx}{}
\newcommand{\intr}{}
\newcommand{\kla}{}
\newcommand{\klam}{}
\newcommand{\klammer}{}
\newcommand{\oper}{}
\newcommand{\tecxt}{}
\newcommand{\Bra}{}
\newcommand{\Ket}{}
\renewcommand{\i}{\text{i}}
\newcommand{\e}{\text{e}}
\newcommand*{\Fourier}{\textsc{Fourier}\ }
\newcommand*{\F}{\mathcal{F}}
\newcommand*{\sinc}{\text{sinc\,}}
\newcommand{\conv}{\mathop{\scalebox{3.5}{\raisebox{-0.38ex}{\makebox[0.5\width][c]{$\ast$}}}}\limits}

\newcommand*{\titel}{Various Problems from Experimental Physics I}
\newcommand*{\autor}{Joachim Schwardt}

\hypersetup{pdfauthor={\autor}, pdftitle={\titel}}

%\titlehead{titlehead}
\subject{Experimental Physics}
\title{\titel}
\author{\autor}
\date{\today}

\begin{document}

%\begin{titlepage}
\maketitle  % Titel setzen
%\tableofcontents  % Inhaltsverzeichnis setzen
%\end{titlepage}

\section{Elevated angled throw}
Let the units of measurement be such that $g=1=v_0$ and consider the elevated angled throw, i.e. $v_{x0} = \cos\alpha$, $v_{y0} = \sin\alpha$ and  $y(t_0) = h$.
Then we find
\begin{align}
x(t) &= t\cos\alpha,\\
y(t) &= -\frac{t^2}{2} + t\sin\alpha + h.
\end{align}
The curve $y(x)$ is thus given by
\begin{align}
y(x) &= -\frac{x^2}{2\cos^2\alpha} + \frac{\sin\alpha}{\cos\alpha} t + h.
\end{align}
The throwing distance is determined by $y(x_g) = 0$ which is a quadratic equation where the positive solution is given by
%%% WOLFRAMALPHA :: Following commands are both understood and yield the correct result :: 
%%% root (\sin\alpha\cos\alpha + \sqrt{\sin^2(\alpha)\cos^2(\alpha) + 2pi\cos^2(\alpha)})'
%%% root (\sqrt{1 + pi - u} ( \sqrt{u - pi} + \sqrt{u + pi}))'
\begin{align}
x_g &= \sin\alpha\cos\alpha + \sqrt{\sin^2\alpha\cos^2\alpha + 2h\cos^2\alpha} = \\
&= \cos\alpha \kla\left( \sin\alpha + \sqrt{\sin^2\alpha + 2h} \right).
\end{align}
Now letting $z := \sin\alpha$ yields
\begin{align}
x_g(z) &= \sqrt{1 - z^2} \kla\left( z + \sqrt{z^2 + 2h} \right).
\end{align}
It is useful to introduce yet another variable $u := z^2 + h$ 
\begin{align}
x_g(u) &= \sqrt{1 + h - u} \kla\left( \sqrt{u - h} + \sqrt{u + h} \right).
\end{align}
Differentiating with respect to $u$ gives
\begin{align}
x_g'(u) &= -\frac{\sqrt{u - h} + \sqrt{u + h}}{2\sqrt{1 + h - u}} + \sqrt{1 + h - u} \kla\left( \frac{1}{2\sqrt{u - h}} + \frac{1}{2\sqrt{u + h}} \right).
\end{align}
To find the optimal throwing angle we have to differentiate with respect to $\alpha$. 
Due to our transformations this condition becomes
\begin{align}
0 &= \frac{\tecxt\text{d} x_g(u(z(\alpha)))}{\tecxt\text{d}\alpha} = \frac{\tecxt\text{d}x(u)}{\tecxt\text{d}u} \frac{\tecxt\text{d}u(z)}{\tecxt\text{d}z} \frac{\tecxt\text{d}z(\alpha)}{\tecxt\text{d}\alpha} = x_g'(u) 2z \cos\alpha.
\end{align}
Since we assume $\alpha \in (0, \frac{\pi}{2})$ this is equivalent to $x_g'(u) = 0$ or
\begin{align}
\sqrt{u - h} + \sqrt{u + h} &= \kla\left( 1 + h - u \right) \left( \frac{1}{\sqrt{u - h}} + \frac{1}{\sqrt{u + h}} \right) = \\
&= (1 + h - u) \frac{\sqrt{u - h} + \sqrt{u + h}}{\sqrt{(u - h)(u + h)}}.
\end{align}
Dividing through by $\sqrt{u + h} + \sqrt{u - h}$ reduces this to
\begin{align}
1 &= \frac{1 + h - u}{\sqrt{u^2 - h^2}}
\end{align}
which is readily transformed to 
\begin{align}
u^2 - h^2 &= (1 + h)^2 + u^2 - 2(1+h)u,\\
\Rightarrow\ \ \ 0 &= 2h^2 + 2h + 1 - 2(1+h)u.
\end{align}
This linear equation has the straight-forward solution
\begin{align}
u &= \frac{2h^2 + 2h + 1}{2 + 2h}.
\end{align}
Thus we find $z=\sqrt{u - h} = \sqrt{\frac{2h^2 + 2h + 1}{2 + 2h} - h}$ and finally the optimal angle
\begin{align}
\alpha &= \arcsin \kla\left( \frac{1}{\sqrt{2 + 2h}} \right).
\end{align}
We may transform back to SI-units by $h\rightarrow h \frac{g}{v_0^2}$ giving
\begin{align}
\alpha  &= \arcsin \kla\left( \frac{v_0}{\sqrt{2v_0^2 + 2gh}} \right).
\end{align}
Using the identity $\tan^2\alpha = \frac{1}{\cos^2\alpha} - 1 = \frac{1}{1 - \sin^2\alpha} - 1$ we may also write this as
\begin{align}
\alpha &= \arctan \kla\left( \frac{v_0}{\sqrt{v_0^2 + 2gh}} \right) \equiv \arctan \kla\left( \frac{v_0}{v_g} \right).
\end{align}
Here we used that the energy of the projectile hitting the ground is $\frac{m}{2}v_0^2 + mgh = \frac{m}{2} v_g^2$, i.e. $v_g = \sqrt{v_0^2 + 2gh}$.
This suggests that there might be a more elegant, presumably geometric approach to reach the final answer.


%\begin{thebibliography}{9}
%\bibitem{example} description, \url{https://www.exmapleurl}, day.\,month~2021.
%\end{thebibliography}

\end{document}